\section{Discussion and bottleneck analysis}

This is where the project earns its title. A roofline plot on its own only shows
that a kernel is slow; it does not say why. The Nsight Compute counters do, and
each kernel's position below is tied to the specific counter that explains it
rather than to a generic claim about what an optimization usually achieves. The
prose in this section is written against the collected counters once the profiling
pass has run on this machine; until then the structure below records exactly
which counter is expected to explain each rung, so the analysis is a fill in of
measured evidence and not a rewrite.

\subsection{The streaming kernels against the bandwidth ceiling}

SAXPY and the reduction should sit close under the measured bandwidth ceiling.
The counter that settles it is the DRAM throughput as a fraction of peak: if
SAXPY reaches most of the measured copy bandwidth, the harness is sound and the
memory bound edge is anchored. Any large shortfall points at launch overhead or
occupancy rather than bandwidth, and the occupancy counter distinguishes them.

\subsection{Transpose: coalescing and bank conflicts}

The naive and tiled transpose is the cleanest counter story in the suite. The
naive kernel's uncoalesced writes show up as a low global store efficiency; the
tiled kernel fixes that and the efficiency counter jumps. The padded tile then
shows up in the shared memory bank conflict counter falling to near zero, where
an unpadded tile would show heavy conflicts on the transposed access. The two
counters, store efficiency and bank conflicts, are exactly the two levers the
tiled transpose pulls.

\subsection{The GEMM ladder: reuse and occupancy}

Up the GEMM ladder, arithmetic intensity and achieved performance should climb
together, and the counters say why at each step. The naive kernel is memory
bound with low intensity. Shared memory tiling raises the intensity by cutting
DRAM traffic, visible as the measured intensity moving right toward the
theoretical value. Register blocking raises reuse further but spends registers,
so the occupancy counter is where its cost shows up: the interesting question is
whether the reuse gain outweighs the occupancy it costs, and the counters answer
it directly. Vectorized loads reduce the memory instruction count, and cuBLAS
sits above the custom kernels as the vendor reference.

\subsection{Thermals and clocks}

Because this is a boost clocked consumer card also driving a display, the NVML
log recorded alongside every run is what shows whether a result was shaped by
temperature or clock drift rather than by kernel design. Where a kernel's timing
variance is large, the joined power and clock samples say whether the card was
throttling or simply boosting unevenly, and that distinction is called out here
wherever the NVML trace makes it.
